% Moderne Paket-Konfiguration für LuaLaTeX
% WICHTIG: Keine inputenc, fontenc oder babel - diese sind für pdfLaTeX!

% --- Sprache ---
\usepackage{polyglossia}
\setmainlanguage[spelling=new]{german}
\setotherlanguage{english}

% --- Kommentar-Blöcke ---
\usepackage{comment}

% --- Mathematische Umgebungen ---
\usepackage{amsmath, amssymb, amsfonts, amsthm, amsbsy}
\usepackage{mathtools}           % Erweitert amsmath
\usepackage{unicode-math}        % Unicode-Mathematik für LuaLaTeX
% ENTFERNT: wasysym, pifont, textcomp (durch unicode-math abgedeckt)

% --- Tabellen und Layout ---
\usepackage{booktabs}            % Schöne Tabellenregeln
\usepackage{multirow}            % Zellen über mehrere Zeilen
\usepackage{array, longtable}
\usepackage{tabularray}          % Moderne Tabellen-Alternative
\usepackage{lscape}              % Seiten im Querformat
\usepackage{pdflscape}           % PDF-optimierte Rotation

% --- Grafiken und Bilder ---
\usepackage{graphicx}            % Bilder einfügen
\usepackage{qrcode}              % QR-Codes generieren
\graphicspath{{media/}}          % Bildpfad global setzen
\usepackage{float}               % Genaues Platzieren
\usepackage{placeins}            % \FloatBarrier
\usepackage{rotating}            % Rotierte Objekte

% --- Bildanordnung (modernisiert) ---
\usepackage{subcaption}          % Modern subcaption (ersetzt subfig)
\usepackage{wrapfig}             % Textumbruch bei Bildern

% --- Symbole und Grafiken ---
\usepackage{fontawesome5}        % Moderne Icon-Fonts (ersetzt marvosym)
\usepackage{tikz}                % Für Diagramme
\usetikzlibrary{positioning, shapes.geometric, arrows.meta}
\usepackage{adjustbox}           % Flexible Bildanpassung

% --- Text und Layout ---
\usepackage{ragged2e}            % Bessere Textausrichtung
\usepackage{setspace}            % Zeilenabstand
\setstretch{1.1}
\usepackage{enumitem}            % Aufzählungen konfigurieren
\usepackage{nowidow}             % Verhindert Schusterjungen/Hurenkinder

% --- Typografie (bereits in Schriftart-settings mit microtype) ---
% microtype wird in Schriftart-settings.tex geladen

% --- Literatur ---
\usepackage[
    backend=biber,
    style=numeric,
    sorting=none,
    sortcites=true,
    giveninits=true,
    maxbibnames=99
]{biblatex}

% --- Glossar ---
\usepackage[acronym,toc,nogroupskip]{glossaries}

% --- PDF und externe Inhalte ---
\usepackage{pdfpages}            % PDFs einfügen
\usepackage{todonotes}           % TODO-Notizen

% --- URLs und Links ---
\usepackage{url}
\usepackage{xurl}                % Bessere URL-Umbrüche

% --- Beschriftungen ---
\usepackage[
    format=plain,
    labelfont=bf,
    justification=centering,
    tableposition=top,
    figureposition=bottom,
    font=small
]{caption}

% --- Farben (für Links, etc.) ---
\usepackage{xcolor}

% --- Listen und Strukturen ---
\usepackage{csquotes}            % Korrekte Anführungszeichen
\MakeOuterQuote{"}               % Deutsche Anführungszeichen
\usepackage{lipsum}              % Platzhalter

% --- Theoremdefinition (modernisiert) ---
\theoremstyle{definition}
\newtheorem{definition}{Definition}[chapter]
\newtheorem{beispiel}[definition]{Beispiel}
\newtheorem{lemma}[definition]{Lemma}
\newtheorem{satz}[definition]{Satz}

% --- Moderne Seitenrotation ---
\newenvironment{rotatepage}{%
    \clearpage
    \begin{landscape}
}{%
    \end{landscape}
    \clearpage
}

% --- Häkchen-Symbol (Unicode-optimiert) ---
\newcommand{\tick}{\textcolor{green!60!black}{✓}}
\newcommand{\cross}{\textcolor{red!60!black}{✗}}

% --- Utility-Befehle ---
\usepackage{xspace}              % Intelligente Leerzeichen
\newcommand{\zB}{z.\,B.\xspace}
\newcommand{\dhh}{d.\,h.\xspace}
\newcommand{\ua}{u.\,a.\xspace}

% --- Debug-Tools (optional) ---
% \usepackage{showframe}         % Zeigt Seitenränder
% \usepackage{lipsum}            % Blindtext

% --- Performance-Optimierungen ---
\usepackage{luatex85}            % Kompatibilität für ältere Pakete

% --- Conditionals für erweiterte Features ---
\usepackage{iftex}
\ifLuaTeX
    % LuaTeX-spezifische Optimierungen
    \usepackage{luacode}
    \usepackage{shellesc}
\fi